\documentclass[10pt,a4paper]{amsart}
\usepackage[margin=19mm]{geometry}
\usepackage{amsmath,amssymb,mathtools}
\usepackage[scr=rsfs,cal=euler]{mathalfa}
\usepackage{xfrac}
\usepackage[unicode,hidelinks,bookmarks=false]{hyperref}
\newcommand{\ie}{i.e.\ }
\newcommand{\cf}{cf.\ }
\newcommand{\resp}{respectively\ }
\newcommand{\Subsection}{\subsection}
\providecommand{\nobreakdash}{\nobreak\hbox{-}}
\setlength{\emergencystretch}{2em}
\multlinegap0pt
\usepackage{amssymb}
\usepackage{mathtools}
\usepackage[shortlabels]{enumitem}
\usepackage{xcolor}
\usepackage{tikz}
\newtheorem{conjecture}{Conjecture}
\newtheorem{theorem}{Theorem}
\newcommand{\IHcomb}{IH_{\mathrm{comb}}}
\newcommand{\IP}{\mathbb{P}}
\newcommand{\IQ}{\mathbb{Q}}
\newcommand{\SigmaFan}{\Sigma}
\title{Combinatorial Cycle Classes in the Intersection Cohomology of Projective Toric Varieties}
\author{Rizwan Jahangir, Daisuke Ishii}
\date{Source: arXiv:2512.06755v2 (CC BY 4.0)}
\begin{document}
\maketitle

\noindent\emph{Source and licence.} Combinatorial Cycle Classes in the Intersection Cohomology of Projective Toric Varieties. Version arXiv:2512.06755v2 (CC BY 4.0), distributed by arXiv under the Creative Commons Attribution 4.0 International licence, http://creativecommons.org/licenses/by/4.0/, as recorded in the arXiv licence metadata for this exact version and shown on the abstract page https://arxiv.org/abs/2512.06755v2. Full licence notice: \texttt{licenses/arxiv-2512.06755-CC-BY-4.0.txt}.\par\smallskip

\noindent\textit{Editorial scope.} Counted unit: the article's theorem theorem (source0.tex lines 217--219). The article's complete original proof of it is delivered with it (source0.tex lines 220--227, one \texttt{proof} environment of 8 lines). No mathematics is written, repaired or re-derived: the statement and the proof are the article's own lines, in the article's own notation, and no retained line is substituted or re-flowed.  Retained uncounted as the source-local context the proof needs: lines 173--180.  Nothing was omitted from any retained span, so the package carries no omission record.  The retained lines cite 4 works, whose entries are transcribed verbatim from the article's own bibliography source in the e-print and delivered with short editorial labels recorded in the item's reference mapping.  No retained line contains a figure environment, an image-inclusion command or drawing code, so no image is delivered and the package declares no asset dependency.  Editorial formatting only, in addition: an AMS \texttt{amsart} preamble that restates the article's own declarations and macros used by the retained lines, namely the environments \texttt{conjecture}, \texttt{theorem} on separate counters, together with the standard packages those lines need.  The retained environments print in this document as Conjecture 1, Theorem 1 in order of appearance, whereas the article prints the same items with the numbering of its own full document.  The complete native source of the article as distributed in the arXiv e-print for this exact version is bundled unchanged as \texttt{sources/190517/source0.tex}.
	\begin{conjecture}[Combinatorial Generation]
		\label{conj:generation}
		For a projective rational fan $\SigmaFan$, the combinatorial cycle classes span the rational even-degree intersection cohomology:
		\[
		\mathrm{span}_{\IQ}\left\{ [V(\tau)]_{\mathrm{comb}} \mid \tau \in \SigmaFan(k) \right\} = \IHcomb^{2k}(\SigmaFan, \IQ)
		\]
		for all $k \ge 0$.
	\end{conjecture}

	\begin{theorem}
		Conjecture~\ref{conj:generation} holds for any projective toric variety of dimension $n \le 2$.
	\end{theorem}

	\begin{proof}
		For $n=1$, $X \cong \IP^1$ is smooth. For $n=2$, let $X$ be a projective toric surface. The relevant degrees are $2k = 0, 2, 4$.
		\begin{itemize}
			\item \emph{Degree 0:} Spanned by the fundamental class $[X]_{\mathrm{comb}}$.
			\item \emph{Degree 2:} By Fieseler \cite[Thm.~1.1]{Fieseler91} and Goresky-MacPherson \cite[6.2]{GM80}, for a normal surface with isolated singularities, $IH^2(X, \IQ) \cong H^2(X, \IQ)$. The ordinary cohomology $H^2$ of a complete toric variety is generated by torus-invariant divisors \cite[Section 3.4]{Fulton}. Thus, the combinatorial cycle classes $[V(\rho)]_{\mathrm{comb}}$ span $\IHcomb^2(\Sigma)$.
			\item \emph{Degree 4:} By Poincar\'e duality, $\IHcomb^4(\Sigma) \cong \IQ$. Since $X$ is projective, it contains smooth points (e.g., vertices of the polytope). For a smooth 2-cone $\sigma$, the local intersection cohomology is trivial, and the cycle class $[V(\sigma)]_{\mathrm{comb}}$ maps to the non-zero generator of $\IHcomb^4$ \cite{BBFK}.
		\end{itemize}
	\end{proof}

\begin{thebibliography}{99}
\bibitem[BBFK]{BBFK} G. Barthel, J.-P. Brasselet, K.-H. Fieseler, and L. Kaup, \textit{Combinatorial intersection cohomology for fans}, Tohoku Math. J. (2) 54 (2002), no. 1, 1--41.
\bibitem[Fiesel]{Fieseler91} K.-H. Fieseler, \textit{Rational intersection cohomology of projective toric varieties}, J. Reine Angew. Math. 413 (1991), 88--98.
\bibitem[Fulton]{Fulton} W. Fulton, \textit{Introduction to Toric Varieties}, Annals of Mathematics Studies, 131. Princeton University Press, Princeton, NJ, 1993.
\bibitem[GM80]{GM80} M. Goresky and R. MacPherson, \textit{Intersection homology theory}, Topology 19 (1980), no. 2, 135--162.
\end{thebibliography}
\end{document}
